\documentclass[11pt,a4paper,twocolumn]{article}

\usepackage[top=1in,bottom=1in,left=0.75in,right=0.75in]{geometry}
\usepackage{amsmath,amssymb,amsthm}
\usepackage{graphicx}
\usepackage{booktabs}
\usepackage{xcolor}
\usepackage{hyperref}
\usepackage{natbib}
\usepackage{microtype}
\usepackage{enumitem}
\usepackage{tikz}
\usetikzlibrary{arrows.meta}
\usepackage{caption}
\usepackage{times}
\usepackage{dblfloatfix}

% ── Give equations room to breathe in narrow columns ─────────────────────────
\setlength{\abovedisplayskip}{9pt}
\setlength{\belowdisplayskip}{9pt}
\setlength{\abovedisplayshortskip}{4pt}
\setlength{\belowdisplayshortskip}{4pt}

\newtheorem{definition}{Definition}
\newtheorem{proposition}{Proposition}

\hypersetup{colorlinks=true,linkcolor=black,citecolor=black,urlcolor=black}

\title{\large\textbf{Nous: A Predictive World Model for Long-Term Agent Memory}}
\author{
  Pranav Singh \\
  Indian Institute of Technology Ropar \\
  \texttt{2023mcb1308@iitrpr.ac.in}
}
\date{\today}

\begin{document}
\maketitle
\thispagestyle{empty}

\begin{abstract}
We present Nous, a novel agent memory architecture grounded in the principle
that knowledge is prediction, not storage. Rather than persisting facts as
database records, vector embeddings, or knowledge-graph triples, Nous
maintains a predictive world model: a collection of categorical probability
distributions, called dimensions, one per entity-attribute pair observed in
conversation. Each incoming observation is scored by its information-theoretic
surprise $\mathcal{S} = -\log_2 P(\text{obs} \mid \mathcal{D})$, and the
distribution is updated via a closed-form Bayesian posterior. The primary
stored artifact is the delta, a record of the shift from prior to posterior
belief, rather than the fact itself. Forgetting emerges naturally as entropy
decay toward the uniform distribution, and identity resolution is handled
through mutual information between entity dimension sets. Evaluated on the
LoCoMo long-term conversational memory benchmark~\citep{maharana2024locomo}
across ten conversations (1,540 questions) using GPT-4o-mini as backbone,
Nous achieves F1 of 63.50 (single-hop), 55.32 (multi-hop), 58.57 (temporal),
and 62.50 (open-domain), outperforming A-MEM~\citep{xu2025amem} by $+24.1$,
$+22.5$, $+27.3$, and $+45.4$ points respectively under identical conditions.
Nous requires no external vector database or graph engine.
\end{abstract}

% ─────────────────────────────────────────────────────────────────────────────
\section{Introduction}
% ─────────────────────────────────────────────────────────────────────────────

A persistent limitation of large language model (LLM) agents is the absence
of a continuously evolving memory. The context window provides short-term
recall within a session, but once a session ends all conversational history
is lost unless externally persisted. This has motivated a growing body of
work on agent memory systems, mechanisms that store, retrieve, and reason
over information accumulated across many sessions.

Existing approaches treat memory as a database of facts. Retrieval-augmented
generation (RAG)~\citep{lewis2020rag} encodes conversation chunks as dense
vectors and retrieves the $k$-nearest at query time. Knowledge-graph systems
such as Zep~\citep{zep2024graphiti} store subject-predicate-object triples.
Episodic systems such as MemGPT~\citep{packer2023memgpt} maintain explicit
core memory strings and flush older records to an archival store. More recent
systems including Mem0~\citep{raghavan2024mem0} and A-MEM~\citep{xu2025amem}
combine graph and vector representations to improve retrieval precision.

Despite their diversity, all these systems share one assumption: memory is a
collection of things that happened. This produces three structural failure
modes. First, conflict blindness: when a user changes jobs or preferences,
contradictory records accumulate with no principled resolution mechanism.
Second, heuristic forgetting: deletion rules are time-based rather than
semantic. Third, retrieval granularity mismatch: a question asking what a
person has painted over the years requires aggregating values from many
sessions, but a fact store returns individual records with no accumulation
mechanism.

We argue these failures share a root cause: treating memory as storage rather
than as a model of the world. The brain does not store a verbatim transcript
of past conversations; it maintains a generative model of its environment and
updates it when observations diverge from predictions
\citep{rao1999predictive,friston2010fep}. Surprise drives learning. Unused
beliefs decay toward uncertainty. Memory is a compressed, continuously refined
probability distribution over the state of the world.

Nous (from the Greek word for mind or intellect) implements this view as a
practical agent memory system. The core data structure is a dimension, a
categorical distribution over the possible values of a single entity-attribute
pair. New observations are scored by information-theoretic surprise and
integrated via Bayesian posterior update. The record of that update, a delta,
is the primary stored artifact, capturing not what was observed but how it
changed the agent's understanding. Forgetting is entropy decay, and conflict
resolution is implicit: contradictory evidence naturally shifts probability
mass away from the old value.

Our contributions are: (1) the predictive world model paradigm for agent
memory, formalised via dimensions, deltas, and surprise-driven Bayesian
updates; (2) a complete open-source implementation requiring no external
database; and (3) controlled experiments on LoCoMo demonstrating consistent
improvement over academic baselines across all four reasoning categories.

% ─────────────────────────────────────────────────────────────────────────────
\section{Background and Related Work}
% ─────────────────────────────────────────────────────────────────────────────

\paragraph{Long-term conversational memory.}
The LoCoMo benchmark~\citep{maharana2024locomo} evaluates agent memory over
conversations spanning up to 35 sessions, 300 to 600 turns, and 9,000 to
16,000 tokens, covering single-hop, multi-hop, temporal, and open-domain
questions. LongMemEval~\citep{wu2025longmemeval} (ICLR 2025) extends this
with harder, audited questions and is increasingly adopted as a complementary
standard.

\paragraph{RAG-based memory.}
RAG~\citep{lewis2020rag} retrieves the top-$k$ embedding-nearest chunks. Its
primary failure is the inability to aggregate information distributed across
many chunks, compounded by the lost-in-the-middle effect~\citep{liu2023lost}
wherein relevant evidence buried in long contexts is ignored.

\paragraph{MemGPT.}
MemGPT~\citep{packer2023memgpt} partitions memory into in-context, recall,
and archival tiers managed via function calls. It is effective for
single-session recall but lacks conflict-resolution semantics. On LoCoMo with
GPT-4o-mini it achieves F1 of 17.29 (single-hop) and 13.18
(temporal)~\citep{xu2025amem}.

\paragraph{A-MEM.}
A-MEM~\citep{xu2025amem} enriches memories with Zettelkasten-style links,
forming a graph of interconnected notes with LLM-generated attributes. It is
the strongest peer-reviewed academic baseline on LoCoMo, reporting F1 of
39.41, 32.86, 31.23, and 17.10 on single-hop, multi-hop, temporal, and
open-domain with GPT-4o-mini.

\paragraph{Zep and Mem0.}
Zep~\citep{zep2024graphiti} builds a temporal knowledge graph; Mem0
\citep{raghavan2024mem0} maintains a hybrid vector and graph store. Both
report strong LoCoMo scores but use GPT-4o with a lenient LLM-as-judge.
A community audit found this judge accepts semantically incorrect answers
roughly 63 percent of the time under adversarial conditions, and approximately
6.4 percent of LoCoMo ground truths contain errors~\citep{locomo2025audit}.
We therefore compare Nous only against systems using the same token-F1 metric
and GPT-4o-mini backbone.

\paragraph{Predictive coding and the Bayesian brain.}
Nous draws theoretically on predictive coding~\citep{rao1999predictive}, which
models the brain as a generative model minimising prediction error, and on
the Free Energy Principle~\citep{friston2010fep}, which unifies perception and
learning under surprise minimisation. The Bayesian brain
hypothesis~\citep{knill2004bayesian} grounds this computationally. Shannon's
information measure~\citep{shannon1948} supplies the definition
$\mathcal{S}(x) = -\log_2 P(x)$ bits, and Cover and Thomas~\citep{cover2006}
supply the entropy and mutual information foundations.

% ─────────────────────────────────────────────────────────────────────────────
\section{The Nous Architecture}
\label{sec:arch}
% ─────────────────────────────────────────────────────────────────────────────

\subsection{Core Abstractions}

\begin{definition}[Dimension]
A dimension $\mathcal{D}_{e,a}$ for entity $e$ and attribute $a$ is a pair
$(V,\, p)$ where $V = \{v_1, \ldots, v_n\}$ is a finite vocabulary and
$p : V \to [0,1]$ satisfies $\sum_{v} p(v) = 1$.
\end{definition}

\medskip

A dimension encodes what the agent believes about one aspect of one entity.
For instance, $\mathcal{D}_{\text{Pranav},\,\text{employer}}$ might be
$\{(\text{Google},\, 0.08),\; (\text{Sarvam AI},\, 0.87),\;
(\text{unknown},\, 0.05)\}$. The mode $\hat{v} = \arg\max_v p(v)$ is the
current best belief.

\begin{definition}[Delta]
A delta $\Delta$ is a tuple
$\langle e,\; a,\; p_{\mathrm{prior}},\; p_{\mathrm{post}},\;
\mathcal{S},\; \tau,\; \omega \rangle$,
where $p_{\mathrm{prior}}$ and $p_{\mathrm{post}}$ are the distributions
before and after the update, $\mathcal{S}$ is surprise in bits, $\tau$ is
the timestamp, and $\omega$ is the supporting evidence text.
\end{definition}

\medskip

The delta is the primary stored artifact of Nous. It records not what happened
but how the agent's understanding changed. The full delta sequence for a
dimension is an auditable history of every belief revision.

\subsection{Information-Theoretic Surprise}

Given dimension $\mathcal{D}_{e,a} = (V, p)$ and observed value obs, the
surprise is:

\begin{equation}
  \mathcal{S}(\text{obs} \mid \mathcal{D}_{e,a})
    = -\log_2\, p(\text{obs})
  \quad \text{bits.}
  \label{eq:surprise}
\end{equation}

If obs $\notin V$, the vocabulary is extended with $p(\text{obs}) = \varepsilon$
and masses renormalised, yielding near-maximum surprise for novel values. An
observation confirming the current belief causes negligible change; one that
contradicts it forces a large revision, mirroring the role of prediction error
in biological memory consolidation~\citep{friston2010fep}.

\subsection{Bayesian Update}

Nous computes the posterior via Bayes' rule~\citep{bayes1763}. The likelihood
is:

\begin{equation}
  \mathcal{L}(\text{obs} \mid v) =
  \begin{cases}
    1 - \varepsilon & v = \text{obs}, \\[4pt]
    \dfrac{\varepsilon}{\,|V|-1\,} & \text{otherwise,}
  \end{cases}
  \label{eq:likelihood}
\end{equation}

and the posterior is:

\begin{equation}
  p_{\mathrm{post}}(v) =
  \frac{\mathcal{L}(\text{obs} \mid v)\; p_{\mathrm{prior}}(v)}
       {\displaystyle\sum_{v' \in V}
        \mathcal{L}(\text{obs} \mid v')\; p_{\mathrm{prior}}(v')}.
  \label{eq:bayes}
\end{equation}

This update is closed-form, runs in $O(|V|)$ time, and requires no gradient
computation. Multiple contradictory observations shift probability mass
naturally without explicit conflict detection.

\begin{proposition}
Under Eq.~\eqref{eq:likelihood}, the posterior mode converges to obs after a
single observation whenever
$p_{\mathrm{prior}}(\text{obs}) > \varepsilon / (1 - \varepsilon)$, which
holds for all realistic vocabulary sizes with $\varepsilon = 0.01$.
\end{proposition}

\subsection{Entropy Decay}

A dimension not updated for time $\Delta t$ grows more uncertain via
exponential mixing toward the uniform distribution:

\begin{equation}
  p_t(v) = \lambda^{\Delta t}\; p_0(v)
          + \bigl(1 - \lambda^{\Delta t}\bigr)\; \mathcal{U}(v),
  \label{eq:decay}
\end{equation}

where $\mathcal{U}(v) = 1/|V|$ and $\lambda \in (0,1)$ is a per-dimension
retention factor. As $\Delta t \to \infty$, $p_t \to \mathcal{U}$, encoding
complete uncertainty. The entropy satisfies:

\begin{equation}
  H(p_t) \;=\; -\sum_{v} p_t(v) \log_2 p_t(v) \;\geq\; \lambda^{\Delta t}\, H(p_0),
  \label{eq:entropy}
\end{equation}

guaranteeing the agent grows less certain as time passes without
reinforcement, a property heuristic deletion rules cannot provide.

\subsection{Identity Resolution}

When two entity mentions may refer to the same person, Nous computes the
symmetrised KL divergence between their shared dimensions:

\begin{equation}
  D(e_1, e_2) = \sum_{a \in \mathcal{A}_{e_1} \cap \mathcal{A}_{e_2}}
  \!\mathrm{KL}(p_{e_1,a} \| p_{e_2,a})
  + \mathrm{KL}(p_{e_2,a} \| p_{e_1,a}).
  \label{eq:kl}
\end{equation}

Entities with low divergence across shared attributes are candidates for
merging via a weighted mixture of their posteriors.

\subsection{Observed-Values Aggregation}
\label{sec:obsvals}

Bayesian posteriors are unimodal: after many updates, mass concentrates on
one value, discarding historical alternatives. For inherently multi-valued
attributes such as activities or places visited, this is incorrect. Nous
addresses this via delta-log aggregation: at query time, every value that
appeared with surprise above threshold $\theta_{\min}$ is included in the
context fact line:

\begin{equation}
  V^*_{e,a} = \bigl\{
    \Delta.\text{obs} \;:\;
    \Delta \in \mathcal{L}_{e,a},\;
    \Delta.\mathcal{S} > \theta_{\min}
  \bigr\}.
  \label{eq:obsvals}
\end{equation}

This preserves the Bayesian world model for the current best belief while
making the full historical record available for aggregation queries.

\subsection{Pipelines}

Figure~\ref{fig:arch} illustrates both pipelines. During ingestion, an LLM
extractor parses each session into $(e,a,v)$ triples. For each triple, the
current dimension is retrieved, the surprise computed via Eq.~\ref{eq:surprise},
the posterior computed via Eq.~\ref{eq:bayes}, and a delta appended to the
log. A compressor prunes dimensions whose entropy exceeds a ceiling threshold.

During querying, named entities in the question seed a breadth-first search
over the entity graph (up to two hops). For each retrieved entity, its
dimensions are formatted as fact lines using Eq.~\ref{eq:obsvals}. The top-$k$
delta records with the highest BM25 similarity to the question are appended as
evidence lines, and the assembled context is passed to a category-specific
prompt and then to the answering LLM.

% ── Full-width figure (explicit coordinates to avoid overlap) ────────────────
\begin{figure*}[t]
\centering
\begin{tikzpicture}[
  font=\small\sffamily,
  box/.style={
    draw=black!50, rounded corners=3pt,
    minimum width=1.9cm, minimum height=0.6cm,
    align=center, fill=white, line width=0.55pt,
    inner sep=4pt
  },
  wbox/.style={
    draw=black!50, rounded corners=3pt,
    minimum width=2.2cm, minimum height=0.6cm,
    align=center, fill=gray!8, line width=0.55pt,
    inner sep=4pt
  },
  arr/.style ={-{Stealth[length=4pt,width=3pt]},
               line width=0.65pt, color=black!55},
  darr/.style={-{Stealth[length=4pt,width=3pt]},
               line width=0.65pt, color=black!40, dashed},
  lbl/.style ={font=\footnotesize\bfseries\sffamily, color=black!50}
]

%% ── Ingestion row (y = 1.5) ──────────────────────────────────────────────
\node[box] (A) at ( 0.0, 1.5) {Session};
\node[box] (B) at ( 3.0, 1.5) {Extract \\ $(e,a,v)$};
\node[box] (C) at ( 6.0, 1.5) {Surprise};
\node[box] (D) at ( 9.0, 1.5) {Update};
\node[box] (E) at (12.0, 1.5) {Delta Log};

\draw[arr] (A)--(B);
\draw[arr] (B)--(C);
\draw[arr] (C)--(D);
\draw[arr] (D)--(E);

%% ── World model (centred, y = 0) ─────────────────────────────────────────
\node[wbox] (W) at (6.5, 0.0) {World Model \\ (Dimensions)};

%% arrows world model <-> ingestion
\draw[arr]  (D.south)  -- ++(0,-0.4) -| (W.north east);
\draw[darr] (W.north)  -- ++(0, 0.4) -| (C.south);

%% ── Query row (y = -1.5) ─────────────────────────────────────────────────
\node[box] (P) at ( 0.0,-1.5) {Question};
\node[box] (Q) at ( 3.0,-1.5) {Entity BFS};
\node[box] (R) at ( 6.0,-1.5) {Profile \\ Assembly};
\node[box] (S) at ( 9.0,-1.5) {Router};
\node[box] (T) at (12.0,-1.5) {LLM Answer};

\draw[arr] (P)--(Q);
\draw[arr] (Q)--(R);
\draw[arr] (R)--(S);
\draw[arr] (S)--(T);

%% world model -> profile assembly
\draw[arr] (W.south) -- ++(0,-0.4) -| (R.north);

%% ── Row labels ───────────────────────────────────────────────────────────
\node[lbl, left=0.3cm of A] {Ingestion};
\node[lbl, left=0.3cm of P] {Query};

\end{tikzpicture}
\caption{Ingestion (top) and query (bottom) pipelines. The world model, a
live store of dimensions, is updated by every Bayesian posterior (solid arrow
from Update) and supplies the current dimension state during surprise
computation (dashed arrow). At query time, the world model feeds the Profile
Assembly step.}
\label{fig:arch}
\end{figure*}

% ─────────────────────────────────────────────────────────────────────────────
\section{Experimental Setup}
% ─────────────────────────────────────────────────────────────────────────────

\subsection{Benchmark}
We evaluate on LoCoMo~\citep{maharana2024locomo}: ten dyadic conversations
spanning 19 to 32 sessions each, totalling 1,540 questions across single-hop
(841), multi-hop (282), temporal (321), and open-domain (96). The adversarial
category is excluded following the canonical protocol.

\subsection{Metrics}
We report token-level F1 and BLEU-1, both strict surface-form measures that
do not reward semantically correct but differently-phrased answers. We
additionally report context recall, the fraction of ground-truth answer tokens
present anywhere in the assembled context, as a retrieval diagnostic, and a
four-way failure bucket classification (good, answer miss, unknown, retrieval
miss) assigned by a GPT-4o-mini judge.

\subsection{Baselines}
All Nous experiments use GPT-4o-mini for extraction, answer generation, and
judging, matching the backbone in \citet{xu2025amem}. We compare against: (i)
a no-memory baseline (direct QA, no retrieved context); (ii) MemGPT
\citep{packer2023memgpt}; and (iii) A-MEM~\citep{xu2025amem}. Baseline scores
are taken from~\citet{xu2025amem}. We do not compare numerically against Zep
or Mem0 due to incompatible backbone and evaluation judge.

% ─────────────────────────────────────────────────────────────────────────────
\section{Results}
\label{sec:results}
% ─────────────────────────────────────────────────────────────────────────────

Table~\ref{tab:main} presents the main results. Nous achieves the highest F1
on all four categories, outperforming A-MEM by between 8.7 and 45.4 points.

\begin{table}[t]
\centering
\caption{Token F1 on LoCoMo (10 conversations, 1,540 questions). All systems
use GPT-4o-mini. Baselines from \citet{xu2025amem}. Best per column in bold.}
\label{tab:main}
\renewcommand{\arraystretch}{1.25}
\setlength{\tabcolsep}{4pt}
\begin{tabular}{lcccc}
\toprule
System & SH & MH & T & OD \\
\midrule
No memory    & 40.36 & 25.02 & 18.41 & 12.04 \\
MemGPT       & 17.29 & 30.36 & 13.18 & 12.24 \\
A-MEM        & 39.41 & 32.86 & 31.23 & 17.10 \\
\midrule
\textbf{Nous}& \textbf{63.50} & \textbf{55.32} & \textbf{58.57} & \textbf{62.50} \\
$\Delta$     & $+24.1$ & $+22.5$ & $+27.3$ & $+45.4$ \\
\bottomrule
\end{tabular}
\end{table}

SH = single-hop, MH = multi-hop, T = temporal, OD = open-domain.

Table~\ref{tab:bleu} reports BLEU-1 and context recall.

\begin{table}[t]
\centering
\caption{Extended results for Nous. Context recall measures retrieval quality
independently of answer quality.}
\label{tab:bleu}
\renewcommand{\arraystretch}{1.25}
\setlength{\tabcolsep}{4pt}
\begin{tabular}{lcccc}
\toprule
Category & F1 & BLEU-1 & Ctx Rec & $n$ \\
\midrule
Single-hop  & 63.50 & 45.60 & 86.07 & 841 \\
Multi-hop   & 55.32 & 23.90 & 75.35 & 282 \\
Temporal    & 58.57 & 56.04 & 61.82 & 321 \\
Open-domain & 62.50 & 10.06 & 47.50 &  96 \\
\midrule
Macro avg   & 59.97 & 33.90 & 67.69 & 1540 \\
\bottomrule
\end{tabular}
\end{table}

The large BLEU-1/F1 gap on open-domain (10.06 vs 62.50) reflects that correct
answers to personality and habit questions can be phrased in many ways, and
BLEU-1 penalises valid paraphrases. On temporal questions BLEU-1 (56.04) and
F1 (58.57) are nearly equal because temporal answers are precise dates or
durations with minimal phrasing variation. The multi-hop gap (BLEU-1 23.90
vs F1 55.32) arises because Nous generates exhaustive lists covering all
ground-truth items plus additional values, raising recall while lowering
unigram precision.

% ─────────────────────────────────────────────────────────────────────────────
\section{Analysis}
\label{sec:analysis}
% ─────────────────────────────────────────────────────────────────────────────

\subsection{Failure Bucket Breakdown}

Table~\ref{tab:buckets} shows the four-way failure distribution. Good denotes
a correct answer; answer miss means context was present but the LLM answered
incorrectly; unknown means the LLM declined to answer despite relevant context
being present; retrieval miss means relevant information was absent from the
assembled context.

\begin{table}[h!]
\centering
\small
\caption{Failure buckets for Nous across all categories.}
\label{tab:buckets}
\renewcommand{\arraystretch}{1.2}
\setlength{\tabcolsep}{3pt}
\begin{tabular}{lrrrr}
\toprule
 & Good & Ans.miss & Unknown & Ret.miss \\
\midrule
SH  & 529 (63\%) & 129 (15\%) & 154 (18\%) & 29 (3\%) \\
MH  &  89 (32\%) &  92 (33\%) &  18 (6\%)  & 24 (9\%) \\
T   & 175 (55\%) &  65 (20\%) &  62 (19\%) & 19 (6\%) \\
OD  &  40 (42\%) &  21 (22\%) &   0 (0\%)  & 35 (37\%) \\
\bottomrule
\end{tabular}
\end{table}

For single-hop, the dominant failure is unknown answer (18\%) despite a
context recall of 86.07\%. Of 154 such failures, 92 occur when context recall
is at least 0.7, meaning the information is present but the model declines to
commit. This is a backbone calibration issue rather than an architectural one.

For multi-hop, good (32\%) and answer miss (33\%) are nearly equal at adequate
context recall (75.35\%), indicating the model struggles to synthesise complete
lists from distributed fact lines. Retrieval miss (9\%) contributes
additionally, as some questions require linking entities beyond the two-hop
search.

Temporal suffers from both low context recall (61.82\%, the lowest of all
categories) and high unknown answer (19\%). Temporal questions often depend on
exact session dates embedded in evidence text rather than extracted as dimension
values, a representational gap addressed in Section~\ref{sec:limitations}.

Open-domain has the highest retrieval miss rate (37\%) and lowest context
recall (47.50\%). Questions about personality and long-term habits do not map
cleanly to entity-attribute dimensions, and improving retrieval is the primary
path forward.

\subsection{Context Recall as a Diagnostic}

Context recall cleanly separates retrieval failures from answering failures.
For single-hop, the 22.6-point gap between context recall (86.07\%) and F1
(63.50\%) identifies the answering step as the binding constraint. For
open-domain, context recall (47.50\%) is itself the ceiling, confirming that
architectural retrieval improvements will directly translate to F1 gains.

% ─────────────────────────────────────────────────────────────────────────────
\section{Limitations and Future Work}
\label{sec:limitations}
% ─────────────────────────────────────────────────────────────────────────────

We do not include controlled ablations isolating the contribution of individual
components such as Bayesian update versus flat storage, entropy decay versus
no forgetting, or delta aggregation versus current-state retrieval. Ablations
are the primary planned extension; the bucket analysis provides directional
evidence for each component's contribution.

We evaluate exclusively on LoCoMo. LongMemEval~\citep{wu2025longmemeval}
provides a harder, more carefully audited benchmark with fewer ground-truth
errors, and evaluation on it is ongoing. Community audits have identified
approximately 6.4 percent erroneous ground-truth answers in LoCoMo and
LLM-judge calibration issues~\citep{locomo2025audit}; our strict token F1
evaluation is therefore conservative.

The categorical dimension is unimodal, requiring delta-log aggregation as a
workaround for set-valued attributes. A set dimension maintaining an
independent Bernoulli distribution per element is the highest-priority planned
architectural extension. Making event timestamps explicit dimension attributes
would additionally improve temporal context recall (currently 61.82\%). All
results use GPT-4o-mini; confirming generalisation across backbone models such
as Llama-3 and Mistral is necessary for a complete picture.

% ─────────────────────────────────────────────────────────────────────────────
\section{Conclusion}
% ─────────────────────────────────────────────────────────────────────────────

We have presented Nous, an agent memory system that treats knowledge as a
predictive world model rather than a database of facts. By representing each
entity-attribute pair as a categorical probability distribution and updating
it via information-theoretic surprise and closed-form Bayesian updates, Nous
handles contradiction, forgetting, and multi-value aggregation in a unified
framework with roots in cognitive science and information theory.

On the LoCoMo benchmark under controlled GPT-4o-mini conditions, Nous
outperforms A-MEM by 24.1, 22.5, 27.3, and 45.4 F1 points on single-hop,
multi-hop, temporal, and open-domain questions respectively. The 86 percent
single-hop context recall confirms that the Bayesian retrieval pipeline
reliably surfaces relevant information; the primary remaining bottleneck is
LLM answer extraction. The results suggest that probabilistic world models are
a competitive and theoretically principled alternative to graph and vector
memory systems. All code and the evaluation harness are released publicly.

\clearpage
\bibliography{nous}
\bibliographystyle{plainnat}

\end{document}
